\subsection{本质上积分特有的一条性质}

下述结果在后面将经常用到。在命题叙述中，不能把本质上积分换成通常的上积分（参看习题 4）。

命题 11. --- 设 \((\lambda_{\alpha})_{\alpha\in\mathrm{A}}\) 是 T 上的正测度族，对关系 \(\leqslant\) 有向，且有上确界 \(\lambda\) （在 \(\mathcal{M}(\mathrm{T})\) 中）。则对每个数值函数 \(f\geqslant0\) ，

\[
\lambda^ {\bullet} (f) = \sup _ {\alpha \in A} \lambda_ {\alpha} ^ {\bullet} (f).\tag{5}
\]

当 \(f\) 属于 \(\mathcal{K}(\mathrm{T})\) 时，此关系式就化为 \(\mathcal{M}(\mathrm{T})\) 中有向集上确界的定义（第二章，§2 第 2 号引理）。其次设 \(f \leqslant g\) 对某个函数 \(g \in \mathcal{K}_+\) 成立（换言之，\(f\) 有界且在一个紧集 \(\mathbf{K}\) 外为零）；设 \(\alpha\) 是使 \(\lambda_{\alpha}(g) \geqslant \lambda(g) - \varepsilon\) 的一个指标，其中 \(\varepsilon\) 是一个数 \(>0\) ；由于测度 \(\nu = \lambda - \lambda_{\alpha}\) 是正的，我们有 \(\nu^{*}(f) \leqslant \nu(g) \leqslant \varepsilon\) ，即 \(\lambda_{\alpha}^{*}(f) \geqslant \lambda^{*}(f) - \varepsilon\) （第四章，§1 第 3 号命题 15）。于是（由于 \(\varepsilon\) 任意）(5) 的右端 \(\geqslant\) 左端；反向不等式是明显的，故对所考虑的特殊情形 (5) 得证。其次，设 \(f\) 在 \(\mathbf{K}\) 外为零但不必有界，并令 \(f_n = \inf(f, n)\) （对每个整数 \(n\) ）。则

\[
\lambda^ {\bullet} (f) = \sup _ {n \in \mathbf {N}} \lambda^ {\bullet} (f _ {n}) = \sup _ {n \in \mathbf {N}} \sup _ {\alpha \in \mathrm{A}} \lambda_ {\alpha} ^ {\bullet} (f _ {n}) = \sup _ {\alpha \in \mathrm{A}} \sup _ {n \in \mathbf {N}} \lambda_ {\alpha} ^ {\bullet} (f _ {n}) = \sup _ {\alpha \in \mathrm{A}} \lambda_ {\alpha} ^ {\bullet} (f).
\]

最后，对 \(f\) 不加任何限制，用 \(\mathfrak{K}\) 表示 \(T\) 的紧子集所成之集，我们有

\[
\begin{array}{c} \lambda^ {\bullet} (f) = \sup _ {K \in \mathfrak {K}} \lambda^ {\bullet} (f \varphi_ {K}) = \sup _ {K \in \mathfrak {K}} \sup _ {\alpha \in A} \lambda_ {\alpha} ^ {\bullet} (f \varphi_ {K}) \\ = \sup _ {\alpha \in A} \sup _ {K \in \mathfrak {K}} \lambda_ {\alpha} ^ {\bullet} (f \varphi_ {K}) = \sup _ {\alpha \in A} \lambda_ {\alpha} ^ {\bullet} (f). \end{array}
\]

推论 1. --- 为使 T 的子集 N 局部 \(\lambda\) -可忽略，必须且只须 N 是局部 \(\lambda_{\alpha}\) -可忽略的（对每个 \(\alpha \in A\) ）。

推论 2. --- 为使 T 到拓扑空间 G 中的映射 g 是 \(\lambda\) -可测的，必须且只须它是 \(\lambda_{\alpha}\) -可测的（对每个 \(\alpha \in A\) ）。

条件显然是必要的，因为 \(\lambda_{\alpha} \leqslant \lambda\) 对每个 \(\alpha\) 成立（第四章，§1 第 3 号命题 15）。反之，设 \(g\) 是 \(\lambda_{\alpha}\) -可测的（对一切 \(\alpha\) ），用 \(\mathfrak{K}\) 表示使 \(g|_{\mathbf{K}}\) 连续的 T 的紧子集 K 所成之集，并设 L 是一个紧集，使 \(\mathrm{L} \cap \mathrm{K}\) 为 \(\lambda\) -可忽略（对每个 \(\mathrm{K} \in \mathfrak{K}\) ）。由于集 \(\mathfrak{K}\) 是 \(\lambda_{\alpha}\) -稠密的，L 是 \(\lambda_{\alpha}\) -可忽略的（对每个 \(\alpha\) ；第四章，§5 第 8 号命题 12），因而是 \(\lambda\) -可忽略的（推论 1）。由此得出 \(\mathfrak{K}\) 是 \(\lambda\) -稠密的，且 \(g\) 是 \(\lambda\) -可测的（第四章，§5 第 10 号命题 15）。

